Smart contracts accept zero-knowledge (ZK) proofs as permission to release
funds, mint credentials, and settle state. A proof certifies only a
statement, yet the contract's action can depend on values the statement never
fixed, such as a recipient taken from calldata. Anyone who observes a valid
proof can then redirect or repeat the effect. We call these \emph{binding
flaws}. Two such failures at the proof-to-contract boundary lost about
\$7.3M in 2025--2026, including \$2.28M drained from Aztec Connect. We define
binding as a relation between a \emph{statement ledger} and an \emph{effect
ledger} under the application's intended policy, and classify its violations
into five patterns that mirror classic signature misuse. In 8,578 public
audit reports we find 87 core binding flaws across 43 project families; 78\%
of those with a severity rating are Critical or High. Existing tools miss
them: Slither, Mythril, circomspect, Picus and generic fuzzing find none of
the audited flaws they could analyze, because the specification lives between
the circuit and the contract. We present \tool{}, which drives an LLM through
an explicit ledger procedure with deterministic analyses, per-proof-system
knowledge and a second look at every verdict, and attaches to each warning
the attack tests it ran on a local chain. A first version found 6 of 22
audited test-set flaws; its misses were errors of judgment, not of reading.
\DEMO{The final version, frozen before any held-out run, finds 10 of 14
held-out audited flaws (4 for the same LLM without the procedure) and 31 of
37 new injected flaws, and it uncovered three previously unknown binding
flaws, each confirmed with a real proof on a local chain and reported to the
maintainers.}
